% tab_model_comparison.tex
\documentclass[tikz, border=2mm]{standalone}
\usepackage{fontspec}
\usepackage[russian]{babel}
\setmainfont{PT Serif}
\setsansfont{PT Sans}
\usetikzlibrary{positioning, arrows.meta, shapes.geometric, shapes}

\begin{document}
\usepackage{tabularx}

\begin{table}[ht]
\centering
\caption{Сравнение подходов к моделированию компьютерных атак}
\label{tab:model_comparison}
\begin{tabularx}{\linewidth}{@{}>{\hsize=0.25\hsize}X
                          >{\hsize=0.25\hsize}X
                          >{\hsize=0.25\hsize}X
                          >{\hsize=0.25\hsize}X@{}}
\toprule
\textbf{Критерий} & \textbf{Деревья атак} & \textbf{Графы атак} & \textbf{Атрибутивные метаграфы} \\
\midrule
Поддержка множественных сущностей & Нет & Частично & Да \\
Семантическая насыщенность & Низкая & Средняя & Высокая \\
Динамическое обновление & Нет & Сложно & Инкрементально \\
Интеграция с MITRE ATT\&CK & Нет & Через маппинг & Нативно \\
Устойчивость к обфускации & Низкая & Средняя & Высокая \\
\bottomrule
\end{tabularx}
\end{table}
\end{document}